\section{Experimental Design}
\label{sec:design}

The evaluation is built from six blocks, arranged so that each one
establishes the noise floor for the next. Every block has pre-registered
validity criteria and a machine-checkable technical-invalid list, written
before the corresponding runs; a run is either a valid outcome or a
technically invalid one, and algorithm failures are never re-run.

\subsection{Simulation Model Validation}

The first block establishes that the simulated plant and its command path
behave reproducibly, which is a precondition rather than a result. It
comprises the per-axis step responses of Table~\ref{tab:stepresponse},
repeated from fresh engine processes, and a closed-loop integration campaign
over three scenarios --- straight transit, a single obstacle ahead, and the
same obstacle with an injected detector dropout --- each repeated from fresh
processes with pre-registered acceptance criteria. Obstacle perception in
this block is the simulator oracle, not the visual detector, so the block
validates physics, control, navigation, and safety plumbing and nothing
about vision.

\subsection{Temporal-Estimation Ablation}

The second block isolates the estimator (RQ4) by holding everything else
fixed. Offline, the deterministic replay harness runs all variants over the
controlled cases D0--D14. In closed loop, the design is
$\{$T0, T1, T2, T3$\} \times \{$E0--E4$\}$ with five fresh-process
repetitions per cell, where E0 is nominal, E1 an early silence shortly after
commitment, E2 a mid-maneuver silence, E3 repeated short silences, and E4 a
large outlier injected shortly after first detection. Success is defined
before the campaign as no collision, minimum clearance above 0.3\,m, exactly
one avoidance entry, no side switch, no abort, and return to the nominal
line. Results are in Sec.~\ref{sec:results}.

\subsection{Planner Benchmark}

The third block compares the two avoidance policies of
Sec.~\ref{sec:planners} on identical perception, using the T2 estimator with
the qualification layer for both. The pre-registered design pairs the
planners by scenario across a long-range set that varies obstacle offset,
obstacle extent, available lateral clearance, mid-maneuver dropout, outlier
injection, initial heading error, and approach speed, plus a pool-scale set
matched to the physical pool geometry. Primary outcomes are collision and
success; secondary outcomes include minimum clearance, maximum lateral
deviation, path length, maneuver time, side switches, thruster effort,
command smoothness, odometry error, and, for the dynamic window, the count of
no-admissible-candidate events. Analysis is by exact collision counts,
success proportions with Wilson intervals, and paired descriptive comparison
of continuous metrics. \textbf{This campaign is in progress and no result
from it appears in this paper:} \PEIGHT

\subsection{Progressive Simulation Calibration}

The fourth block instantiates the ladder of Sec.~\ref{sec:calibration}. Each
level requires an independent measurement campaign against the physical
system --- camera characterization at \SONE{}, a latency budget at \STWO{},
and step-response identification plus sensor-noise characterization at
\STHREE{} --- followed by re-running the simulation conditions of block C at
that level. The measurement procedures are specified; the measurements
themselves require the physical vehicle and the pool. \CALIB

\subsection{Physical Pool Experiments}

The fifth block runs the matched conditions on the physical vehicle. The
pool is rectangular, roughly 8\,m long, shallow, with a V-shaped floor and no
controlled currents; exact dimensions, the bottom profile, obstacle
dimensions and positions, and the start points are to be surveyed by a
written measurement protocol before any run, and are deliberately not
assumed here. Two physical obstacles are used, an anchor and a torpedo.
Because the pool is short, the scenarios are pool-scale by construction:
reduced speed, shorter initial obstacle distance, a compact avoidance
maneuver at approximately constant depth, horizontal avoidance using surge,
sway and yaw, and physical safety boundaries with the avoidance side
constrained never to point into the nearer wall. The long-range trajectories
of block A are therefore \emph{not} representative of the physical
experiment, and the matched simulation conditions use the pool-scale profile
on both sides.

Validity criteria are pre-registered: start pose within a fixed tolerance,
all logs recording, ground-truth marker visible for at least 95\% of the run,
and no operator input after the start signal. Collision, near collision,
successful avoidance, successful recovery, and operator intervention each
have explicit definitions, and technically invalid runs are archived with
their reason and counted in the report rather than silently dropped.
Condition order is randomized within blocks and the operator does not see the
algorithm identity during a run. \POOL

\subsection{External Ground Truth}

The sixth block is the measurement instrument for block E. An Intel
RealSense D435 provides validation-only pose: the inventoried device
sustains $1920{\times}1080$ RGB at a measured 29.81\,fps with host-clock
timestamps, which removes cross-device clock transfer since it shares the
experiment machine. It is to be mounted rigidly near nadir over the test
area, with at least four measured non-collinear perimeter references defining
the pool frame and a high-contrast fiducial board on the vehicle providing
$x$, $y$, and yaw; depth comes from the vehicle's pressure sensor.

Depth measured through the air--water interface is \emph{not} trusted as
ground truth. Refraction makes the pixel-to-pool mapping depth-dependent, so
the planned procedure is a dry intrinsic calibration, extrinsics from the
above-water perimeter markers, and a multi-plane submerged calibration at
several surveyed depths interpolated by the vehicle's measured depth, with an
explicit flat-refraction model as the fallback if residuals require it. The
instrument is usable only after its own error report --- position and yaw
RMSE, 95th percentiles, error versus depth, timestamp uncertainty --- is
documented. No ground-truth accuracy figure is claimed until then. \POOL

Comparable external references exist in this literature and set the rigor
bar~\cite{winkel2023designim,yang2025duvindif}; the distinguishing constraint
here is that the reference is validation-only and never enters estimation,
unlike overhead cameras used as a measurement source inside the
filter~\cite{alineipoiana2024abluerov}.
